% 第二章 系统相关技术概述
\chapter{系统相关技术概述}

本章对"今天吃什么"智能推荐系统所依赖的理论基础与关键技术进行系统阐述。首先介绍推荐系统的基本范式与餐饮推荐的特殊约束，继而分析主流推荐算法的原理与局限；随后分别从前端、后端、数据库、文本向量化及安全认证五个维度，深入剖析各技术的核心机制，并结合本项目的业务特征论证技术选型的合理性，为后续系统设计与实现奠定理论基础。

\section{推荐系统基本原理}

\subsection{推荐系统的定义与目标}

推荐系统（Recommender System）本质上是一种信息过滤机制，其核心任务是在用户集合 $U$ 与物品集合 $I$ 构成的二元空间上，学习一个评分映射函数 $f: U \times I \rightarrow \mathbb{R}$，使得对于特定用户 $u \in U$，系统能够从其尚未交互过的物品子集 $I_u^- = I \setminus I_u^+$ 中筛选出预测评分最高的前 $N$ 个物品，形成个性化推荐列表 $R_u$。形式化表达为：

\begin{equation}
R_u = \mathop{\arg\max}_{i \in I_u^-} f(u, i)
\end{equation}

该定义揭示了推荐系统的两个核心目标：一是准确性，即推荐结果应尽可能贴合用户的真实偏好；二是覆盖率，即系统应能够在海量候选中有效挖掘长尾物品，避免推荐结果过度集中于热门项目。在餐饮推荐场景中，这一目标还需进一步受限于预算、距离、饮食禁忌等现实条件，使得评分函数 $f$ 不再仅仅是偏好强度的度量，而是可执行性的综合评估。项亮在《推荐系统实践》中系统阐述了推荐系统的工程化方法，涵盖从数据采集、特征工程到算法评估的完整流程\super{[16]}。

\subsection{推荐系统的典型架构}
在实际应用中，工业级推荐系统大多采用多阶段级联的架构形式，以此平衡计算效率与推荐精准度。整体架构通常分为召回、粗排、精排和重排四个环节。召回阶段的主要任务是快速筛选候选，从海量物品库中提取出上千条相关候选，常用方式包括协同过滤召回、向量相似度召回和规则标签召回等，这一阶段更注重筛选速度与覆盖范围，允许存在少量精度损失。粗排阶段会对召回结果做初步精简，借助逻辑回归、浅层神经网络等轻量模型快速打分，把候选规模压缩到合适范围，减少后续精排的计算压力。精排是提升推荐效果的关键环节，通过DeepFM、DIN等复杂深度学习模型，对候选菜品进行精细化个性化打分，充分挖掘用户与物品之间的关联特征。最后的重排阶段，会结合业务规则对结果做优化调整，例如控制推荐多样性、平衡菜品新鲜度、避免重复推荐，同时加入硬性约束过滤，最终生成可直接展示给用户的推荐列表。

这种多阶段分层架构的核心思路，是把不同复杂度的计算任务分配到对应阶段，在保证响应速度的前提下提升推荐效果。本系统所采用的“硬性过滤—向量召回—综合重排”策略，正是对该架构的轻量化适配，先通过规则约束保证推荐可用，再利用语义向量提升匹配精度，最后结合多目标因素完成排序优化。

\subsection{餐饮推荐的特殊性}
与电商、新闻等常见推荐场景相比，餐饮推荐具备鲜明的领域特点，这些特点也直接决定了系统的设计方向。首先，餐饮推荐具有较强的约束性，用户的就餐选择会被预算、距离、营业状态、饮食禁忌等条件严格限制，任何超出约束范围的推荐都不具备实际使用价值。其次，餐饮需求具有明显的时效性，用户在早午晚餐、夜宵等不同时段的偏好差异较大，即使是同一时段，也会因身体状态、心情变化产生不同需求，单纯依靠历史行为很难适配这种动态变化。再次，餐饮消费的容错率较低，一旦完成点餐，后续修改或更换的成本较高，因此对推荐结果的准确性要求更高。最后，餐饮推荐需要兼顾多目标平衡，用户通常会同时考虑口味、健康、价格、距离等多个因素，系统需要在这些维度之间做好权衡，而不是只优化单一指标。

结合这些特点，本系统没有采用仅依赖历史行为的协同过滤方案，而是采用静态画像结合即时问答的双通道建模方式，通过主动交互获取当前场景信息，并将硬性约束放在推荐流程的前端，保证最终结果符合实际使用场景。

\section{常用推荐算法分析}
\subsection{基于内容的推荐算法}
基于内容的推荐算法，核心是依靠物品自身的属性特征与用户的历史偏好做匹配。具体实现时，系统会先提取物品的关键特征，构建物品特征向量；同时根据用户的点击、收藏、下单等历史行为，总结出用户的兴趣偏好向量；最后通过计算两种向量的相似度，把相似度较高的物品推荐给用户。在餐饮场景中，菜品的菜系、口味、食材、价格、营养成分等，都可以作为特征参与计算。

这种算法的优势在于，用户建模相互独立，不会受其他用户数据的干扰，新上架的菜品也能立刻参与推荐，不存在物品冷启动的问题。但它也存在明显局限，推荐结果会被用户的历史偏好限制，很难挖掘用户潜在的兴趣点，容易出现推荐内容同质化的问题；同时对于没有任何行为记录的新用户，无法构建偏好模型，仍然会遇到用户冷启动的难题。在本系统中，基于内容推荐的思想主要应用在菜品向量化部分，通过提取菜品的多维属性生成语义向量，为后续的相似度计算提供支撑。周志华在《机器学习》中系统介绍了各类学习算法的原理与适用场景，为本系统的算法选型提供了理论参考\super{[17]}。


\subsection{协同过滤推荐算法}

协同过滤（Collaborative Filtering, CF）是当前应用最为广泛的推荐算法之一，其核心假设是"兴趣相似的用户倾向于偏好相似的物品"。该算法不依赖物品本身的内容特征，而是基于用户与物品之间的交互矩阵（如评分矩阵或隐式反馈矩阵）进行建模。从实现路径上，协同过滤可分为基于用户的协同过滤（User-based CF）与基于物品的协同过滤（Item-based CF）两类。

基于用户的协同过滤首先计算目标用户与其他用户之间的相似度（常用皮尔逊相关系数或余弦相似度），选取相似度最高的 $K$ 个用户作为近邻集合；随后根据近邻用户对候选物品的评分预测目标用户的偏好程度，将预测评分最高的物品予以推荐。基于物品的协同过滤则转换视角，先计算物品之间的共现相似度（即同时被同一用户偏好的概率），再根据目标用户历史喜欢的物品，推荐与其历史物品相似度最高的其他物品。

协同过滤算法能够突破用户既有兴趣边界的限制，通过群体行为发现潜在关联，对长尾物品的挖掘能力较强。但其固有缺陷在于：新用户或新物品因缺乏历史交互数据而难以被有效建模，即冷启动问题；随着用户规模和物品 catalog 的增长，用户-物品交互矩阵会极度稀疏，导致相似度计算失真；此外，算法对数据规模高度敏感，初期数据不足时推荐质量难以保证。考虑到本系统面向的是即时决策场景，用户的历史餐饮记录有限且场景波动大，单纯依赖协同过滤难以适应"当前想吃什么"这一动态需求，因此本系统将其作为辅助思路，而非核心推荐策略。李航在《统计学习方法》中对监督学习的主要模型做了严密的数学推导与算法总结，其中的分类与回归方法为本系统的评分预测提供了方法论基础\super{[18]}。

\subsection{混合推荐策略}

混合推荐并非一种独立的算法，而是将多种推荐算法按照特定策略进行组合，以发挥各自优势、弥补单一算法的不足。常见的组合模式包括加权融合、层叠过滤、特征扩充与元级联等。加权融合模式并行执行多种算法，通过线性加权整合各算法的预测得分；层叠模式则采用串行结构，先以一种算法生成候选集合，再以另一种算法对候选进行精细化筛选或重排；特征扩充模式将一种算法生成的中间特征作为另一种算法的输入，实现信息互补。

本系统采用的推荐策略属于层叠混合模式，具体表现为"硬性过滤---向量召回---综合重排"的三段式流水线。第一阶段基于规则进行硬性过滤，利用预算区间、审核状态及饮食禁忌等约束条件快速剔除不可执行候选，缩小搜索空间；第二阶段基于语义向量进行相似度召回，捕捉用户当前意图与菜品属性之间的语义关联；第三阶段综合语义相似度、距离衰减与历史平滑等因素进行重排序，实现多目标权衡。这种分层架构既保证了约束的严格满足，又通过向量语义提升了推荐的精准度，同时避免了单一算法在复杂场景下的局限性。

\section{前端开发技术}

\subsection{微信小程序技术架构}
微信小程序是腾讯在2017年推出的轻量级应用形式\super{[19]}，它采用逻辑层与渲染层相互分离的双线程架构。其中，逻辑层基于独立的JavaScript引擎JSCore运行，主要负责处理业务逻辑、管理数据状态以及发起网络请求；渲染层则负责页面的视觉渲染，使用WXML描述页面结构，通过WXSS定义页面样式。两层之间无法直接通信，而是借助底层Native提供的JSBridge机制完成异步数据交互，逻辑层通过`setData`方法将数据更新传递到渲染层，从而触发页面重新渲染。

这种双线程隔离的设计带来了两方面突出优势。一方面提升了运行安全性，由于逻辑层不具备完整的DOM操作权限，开发者无法直接操作页面节点，能够有效避免XSS注入等前端安全风险。另一方面优化了渲染性能，WXML和WXSS代码经过编译后由Native层直接解析渲染，减少了传统WebView中频繁的DOM计算与页面重绘，提升了页面流畅度。同时，微信小程序提供了网络请求、本地存储、地理位置、扫码等丰富的原生接口，开发者可以通过统一的API快速调用系统能力，降低功能实现难度。

在组件开发方面，小程序内置了视图容器、基础控件、表单、导航等常用组件，同时支持自定义组件的封装与复用。在样式适配方面，WXSS使用rpx作为自适应单位，以屏幕宽度750rpx为基准进行等比缩放，能够自动适配不同分辨率的移动设备，有效降低了多机型适配的开发成本。

\subsection{微信小程序在本项目中的适用性分析}
本系统最终选用微信小程序作为前端展示载体，主要结合了目标用户场景与实际开发效率两方面因素。从用户使用角度来看，高校学生几乎都在使用微信，小程序无需下载安装、即用即走的特点，大幅降低了用户的使用门槛；同时小程序支持消息订阅和社交分享，便于在用餐高峰时段向用户推送推荐内容，提升系统使用率。从技术实现角度来看，小程序的双线程架构可以稳定支撑问答交互页面的分步操作和推荐结果页面的动态加载，其原生提供的地理位置接口能够直接获取用户经纬度，为后续距离计算和附近美食推荐提供了可靠的数据来源。此外，微信小程序拥有完善的开发者工具，支持代码编写、真机调试和性能监控，能够满足本项目从开发、调试到测试的完整工程需求。


\section{后端开发框架}

\subsection{FastAPI 框架技术原理}

FastAPI 是基于 Python 3.7 及以上版本类型注解构建的现代 Web 框架\super{[20]}，其底层建立在 Starlette ASGI 工具包与 Pydantic 数据验证库之上。ASGI（Asynchronous Server Gateway Interface）是 Python 异步服务器网关接口标准，相较于传统的 WSGI（Web Server Gateway Interface），ASGI 原生支持异步非阻塞 I/O 操作，使得框架能够在单线程内通过事件循环（Event Loop）并发处理大量连接，显著提升了 I/O 密集型场景下的吞吐量。

FastAPI 的核心设计哲学是将 Python 的类型系统（Type Hints）作为一等公民。开发者在路由处理函数的参数中声明数据类型后，FastAPI 自动利用 Pydantic 完成请求体的解析、校验与序列化。Pydantic 通过 Python 的类型注解在运行时构建数据模型，对传入的 JSON 数据进行严格的类型检查与约束验证，一旦数据不符合模型定义，框架会自动返回结构化的 422 错误响应，无需开发者编写冗余的校验代码。这种类型驱动开发模式不仅减少了样板代码，还使得接口的数据契约具有自文档化特性。

在依赖管理方面，FastAPI 提供了基于 \texttt{Depends} 装饰器的依赖注入系统。开发者可以将数据库会话获取、用户身份校验、权限检查等横切关注点封装为依赖函数，通过声明式的方式注入到路由处理函数中。这种设计实现了业务逻辑与基础设施的解耦，使得代码具有更好的可测试性与可复用性。此外，FastAPI 内置对 OpenAPI 与 JSON Schema 的自动生成支持，服务启动后会自动暴露交互式 API 文档（Swagger UI 与 ReDoc），前端开发者可直接在浏览器中查看接口定义、参数结构并发起调试请求，显著提升了前后端协作效率。

\subsection{SQLAlchemy ORM 技术}

SQLAlchemy 是 Python 生态中功能较为完备的对象关系映射（ORM）工具集，其设计理念是在保持关系型数据库完整表达能力的同时，提供高层次的数据操作接口。SQLAlchemy 的 ORM 层采用声明式映射（Declarative Mapping）模式，允许开发者通过定义 Python 类来描述数据库表结构，类的属性对应表的列，类实例对应表中的行记录，从而将数据库操作转化为直观的面向对象操作。

在会话管理方面，SQLAlchemy 通过 Session 对象实现单元工作模式（Unit of Work）。Session 作为对象与数据库之间的持久化上下文，跟踪所有待持久化对象的状态变更，并在事务提交时（\texttt{session.commit()}）将变更批量同步至数据库；若操作过程中发生异常，可通过回滚（\texttt{session.rollback()}）撤销未提交的变更，确保数据一致性。此外，SQLAlchemy 内置连接池（Connection Pool）机制，对数据库连接进行复用与管理，避免了频繁创建连接带来的性能开销。

SQLAlchemy 的另一重要特性是数据库无关性。通过更换连接字符串与方言（Dialect）配置，同一套 ORM 代码可无缝迁移至 MySQL、PostgreSQL、SQLite 等不同数据库后端，这种抽象能力为开发环境与生产环境的灵活部署提供了便利。在本系统中，SQLAlchemy 的声明式基类被用于定义用户、画像、菜品、店铺、历史记录等核心实体模型，并通过外键关系与级联操作实现了实体间的关联建模。

\subsection{后端技术选型分析}

在 Python Web 开发领域，FastAPI、Flask 与 Django 是三种主流框架，三者在设计理念与适用场景上存在显著差异。Flask 作为微框架，核心功能精简，扩展依赖第三方库，学习曲线平缓，适合快速原型开发，但缺乏原生异步支持与自动文档生成能力，大型项目中需自行整合校验、路由及文档工具。Django 作为全栈框架，内置 ORM、管理后台、认证系统等完整功能，适合内容密集型的大型应用，但其同步架构与较重的框架包袱对于轻量级 API 服务而言显得冗余，且学习成本较高。FastAPI 则定位于现代异步 API 框架，原生支持 ASGI 异步协议与类型注解驱动开发，在接口规范性、并发性能及开发效率之间取得了良好平衡。

本系统的后端属于典型的接口驱动型架构，前后端通过 RESTful API 进行数据交互，对自动文档生成、数据校验及异步并发均有较高要求。推荐接口涉及向量计算、数据库查询及地理距离运算等 I/O 密集型操作，FastAPI 的异步处理能力能够有效降低接口延迟；同时，系统在开发阶段需要频繁进行前后端联调，FastAPI 自动生成的 Swagger 文档显著降低了接口沟通成本。综合上述因素，本系统选择 FastAPI 作为后端主框架，并以 SQLAlchemy 作为数据持久化层。


\section{数据库技术}

\subsection{MySQL 数据库技术}
MySQL是目前使用范围最广的开源关系型数据库，它默认搭载的InnoDB存储引擎，可以提供完整的事务支持和系统崩溃后的数据恢复能力，保证数据稳定可靠。InnoDB采用多版本并发控制（MVCC）机制实现高效的并发读写：当事务对数据进行修改时，不会直接覆盖原有记录，而是生成一条新的数据版本，并通过回滚指针维护版本之间的关联；读取数据时，则按照事务的隔离级别读取对应可见版本的快照，从而避免传统锁机制中读写冲突带来的阻塞问题，进一步提升数据库的并发处理性能。

在索引结构上，InnoDB以B+树作为聚簇索引和二级索引的底层实现结构。聚簇索引的叶子节点直接存放完整的行数据，因此基于主键的查询效率极高；二级索引的叶子节点只存储主键值，需要通过回表操作才能查到完整记录。合理设计索引可以明显减少查询的I/O开销，提升数据检索速度。此外，MySQL支持标准SQL语法和丰富的多表关联查询（JOIN），可以高效处理用户、菜品、画像等多表之间的复杂关联操作，还能通过外键约束保证数据表之间的引用完整性。

事务的ACID特性是InnoDB保障数据安全的核心。其中，原子性依靠undo log实现，确保一个事务内的操作要么全部提交，要么全部回滚；持久性依靠redo log实现，即使系统崩溃，也能根据日志恢复已经提交的事务数据；隔离性由MVCC和锁机制共同实现。MySQL提供读未提交、读已提交、可重复读、串行化四种隔离级别，默认使用的可重复读级别，能够在大多数业务场景下兼顾数据一致性和并发效率。


\subsection{数据库选型依据}

本系统涉及用户、用户画像、店铺、菜品、历史记录及交互日志六类核心实体，实体之间存在复杂的多对一、一对多关联关系（如店铺与菜品、用户与历史记录）。关系型数据库在这种强关联、强一致性的场景下具有天然优势，能够通过 JOIN 操作高效完成跨表查询，并通过事务机制保证推荐批次写入与选择确认等复合操作的数据一致性。

相较于文档型数据库 MongoDB，MySQL 在复杂关联查询与事务支持方面更为成熟；相较于嵌入式数据库 SQLite，MySQL 的客户端-服务器架构具备更高的并发处理能力与独立运维空间，适合作为多用户同时访问的后端数据库。因此，本系统选择 MySQL 作为主数据库，并通过 SQLAlchemy ORM 进行对象关系映射，兼顾开发效率与运行性能。

\section{文本向量化技术}

\subsection{从传统文本表示到预训练语言模型}

文本向量化旨在将非结构化的自然语言文本映射为低维稠密向量，使得语义相近的文本在向量空间中距离较近。早期的文本表示方法以词袋模型（Bag of Words, BoW）和 TF-IDF 为代表，这类方法将文本视为词汇的无序集合，通过统计词频构建高维稀疏向量，虽然计算简单，但完全忽略了词语之间的语义关联与上下文信息，无法区分"不喜欢"与"不，喜欢"这类语序敏感表达。

分布式词表示（如 Word2Vec、GloVe）通过神经网络在大规模语料上学习词的低维嵌入，使得语义相近的词在向量空间中聚集，一定程度上缓解了稀疏性与语义鸿沟问题。然而，词向量本质上是静态的，无法处理一词多义现象，且缺乏对句子级别语义的建模能力。

2018 年，Google 提出的 BERT（Bidirectional Encoder Representations from Transformers）模型通过 Transformer 编码器与双向上下文预训练\super{[21]}，彻底改变了文本表示的范式。BERT 采用掩码语言模型（Masked Language Model, MLM）任务，在预训练阶段随机遮蔽输入序列中的部分 token，要求模型根据上下文预测被遮蔽的词，从而迫使模型学习到深层的双向语义表示。相较于之前的单向语言模型，BERT 能够同时利用词语的左侧与右侧上下文信息，生成的上下文相关词向量在语义理解任务上表现优异。

\subsection{Sentence-BERT 模型架构}

尽管 BERT 在词级别和句级别任务上均取得了突破性进展，但其原生输出是 token 级别的隐藏状态，若直接对句子所有 token 的向量取平均或取 [CLS] 向量作为句子表示，在语义相似度任务上的效果并不理想。Reimers 与 Gurevych 提出的 Sentence-BERT（简称 SBERT）通过双塔 Siamese 网络架构解决了这一问题\super{[22]}，专门优化句子级别的语义嵌入质量。

SBERT 的核心结构由两个共享权重的编码器（通常基于 BERT 或 RoBERTa）组成。输入的两个句子分别经过编码器处理，得到各自的句子向量 $u$ 和 $v$；随后通过池化策略（如均值池化或 [CLS] 池化）将变长的 token 序列压缩为固定维度的句子向量。在训练阶段，SBERT 采用对比学习范式：对于语义相似的句子对（正例），优化目标使向量距离尽可能小；对于语义不相似的句子对（负例），使向量距离尽可能大。常用的损失函数包括对比损失（Contrastive Loss）与三元组损失（Triplet Loss），其中三元组损失通过锚定句子（Anchor）、正例句子（Positive）和负例句子（Negative）的相对距离关系进行优化，要求锚定与正例的距离显著小于锚定与负例的距离。

经过有监督的句子对训练后，SBERT 生成的句子向量具有优良的语义对齐特性，可直接通过余弦相似度或欧氏距离度量句子间的语义关联，且计算复杂度远低于 BERT 的交叉编码方式，适合大规模语义检索与召回场景。

\subsection{中文语义向量表示}

本系统采用 \texttt{shibing624/text2vec-base-chinese} 作为句子编码器，该模型是基于中文语料在 SBERT 架构上微调的语义向量模型。模型以中文 BERT 为基座，通过中文语义文本相似度数据集进行有监督训练，能够将中文句子编码为 768 维的稠密向量。在餐饮推荐场景中，系统将菜品的名称、菜系、口味标签、食材及描述等结构化字段拼接为一段语义文本，通过该模型编码为菜品向量；同时，将用户的静态画像与即时问答结果拼接为用户语义文本，编码为用户向量。基于向量的语义召回使得系统能够在没有大量历史交互数据的情况下，通过语义相似度捕捉用户当前意图与菜品属性之间的潜在关联，有效缓解了冷启动问题。

\section{身份认证与安全机制}

\subsection{JWT 无状态认证机制}

JSON Web Token（JWT）是一种基于 JSON 的开放标准（RFC 7519），用于在网络应用间安全地传输认证信息。JWT 由 Header、Payload 和 Signature 三部分组成，各部分分别进行 Base64Url 编码后通过点号连接形成完整令牌。

Header 部分声明令牌类型与签名算法（如 HS256）；Payload 部分承载声明（Claims），包括签发者、主题（通常为用户标识）、签发时间（iat）及过期时间（exp）等；Signature 部分则使用服务端密钥对前两部分进行 HMAC-SHA256 签名，确保令牌内容在传输过程中未被篡改。由于签名依赖于服务端持有的密钥，攻击者无法在不知道密钥的情况下伪造有效令牌。

JWT 的本质是一种无状态（Stateless）认证机制。服务端在验证令牌时无需查询数据库或缓存中的会话状态，仅通过校验签名与过期时间即可确认用户身份，这显著降低了服务端的状态维护成本，提升了系统的水平扩展能力。然而，无状态特性也意味着令牌一旦签发，在过期之前无法被服务端主动撤销，因此需要合理设置令牌有效期，并在敏感操作中结合其他安全策略。

\subsection{本项目安全访问设计}

本系统采用 JWT 实现登录态管理与接口访问控制。用户登录成功后，后端使用 \texttt{python-jose} 库签发 access\_token，Payload 中嵌入用户唯一标识与用户名，令牌有效期设置为 24 小时。前端在后续请求的 HTTP Header 中以 \texttt{Authorization: Bearer <token>} 格式携带令牌。后端通过依赖注入机制在每个受保护接口中统一校验令牌：首先验证签名是否有效，其次检查令牌是否过期，最后查询数据库确认用户存在且账号状态正常。对于已封禁用户，即使持有有效令牌，系统也会返回 HTTP 403 状态码拒绝访问，实现业务层面的权限隔离。

这种设计将认证逻辑从业务代码中剥离，通过声明式依赖注入统一处理，既保证了接口的安全性，又避免了在每个路由中重复编写校验代码，使系统具备良好的可维护性与扩展性。

\section{本章小结}

本章介绍了系统涉及的主要技术。推荐系统一般采用多阶段架构来兼顾效率与精度，本系统将其简化为"硬性过滤—向量召回—综合重排"的三段流程。算法方面，协同过滤虽然应用广泛，但对历史数据依赖较重，不太适合即时餐饮场景；基于内容的推荐可解释性好但容易同质化；因此本系统选择了混合策略。前端选用微信小程序，主要考虑高校学生几乎都在用微信，不需要额外安装，且原生支持定位功能。后端用 FastAPI 搭配 SQLAlchemy，前者能自动生成接口文档、方便前后端联调，后者通过 ORM 屏蔽了数据库差异。数据存储选 MySQL，主要是事务支持和复杂查询能力比较成熟。向量化部分用 sentence-transformers 的本地模型把菜品描述转成向量，避免依赖外部接口。安全方面用 JWT 做无状态认证，配合业务层的账号状态检查。这些技术共同构成了系统的实现基础。
